\documentclass[parskip=full]{scrartcl}
\usepackage{geometry}
\geometry{a4paper, left=2.5cm, right=2cm, top=3cm, bottom=3cm}
\usepackage[utf8]{inputenc} % use utf8 file encoding for TeX sources
\usepackage[T1]{fontenc}    % avoid garbled Unicode text in pdf
\usepackage[german]{babel}  % german hyphenation, quotes, etc
\usepackage{hyperref}       % detailed hyperlink/pdf configuration
\hypersetup{                % ‘texdoc hyperref‘ for options
pdftitle={Echtzeit-Computergrafik in der Spieleentwicklung},%
bookmarks=true,%
}

\usepackage{rotating}
\usepackage{newclude}
\usepackage{typearea}
\newcommand\invisiblesection[1]{%
  \refstepcounter{section}%
  \addcontentsline{toc}{section}{\protect\numberline{\thesection}#1}%
  \sectionmark{#1}}
  

\usepackage{listings}

\usepackage{xcolor}
\usepackage{graphicx}       % provides commands for including figures
\usepackage{csquotes}       % provides \enquote{} macro for "quotes"
\usepackage[nonumberlist]{glossaries}     % provides glossary commands
\usepackage{enumitem}
\usepackage{epsf,epsfig,eepic}
\usepackage[headsepline, footsepline]{scrpage2}
\usepackage{tabularx}


%Design für Methoden-/Attributbeschreibungen  
\usepackage{framed}
\usepackage{color}
\definecolor{gray}{RGB}{200,200,200}
\renewenvironment{leftbar}[1][\hsize]
{%
\def\FrameCommand
{%
{\color{gray}\vrule width 1pt}%     
}%
\MakeFramed{\hsize#1\advance\hsize-\width\FrameRestore}%
}
{\endMakeFramed}

\definecolor{eclipseStrings}{HTML}{B71C1C}
\definecolor{eclipseKeywords}{HTML}{0D47A1}
\definecolor{background}{HTML}{ECEFF1}
\colorlet{numb}{magenta!60!black}

\lstdefinelanguage{json}{
    basicstyle=\normalfont\ttfamily,
    commentstyle=\color{eclipseStrings}, % style of comment
    stringstyle=\color{eclipseKeywords}, % style of strings
    numbers=left,
    numberstyle=\scriptsize,
    stepnumber=1,
    numbersep=8pt,
    showstringspaces=false,
    breaklines=true,
    frame=lines,
    backgroundcolor=\color{background}, %only if you like
    string=[s]{"}{"},
    comment=[l]{:\ "},
    morecomment=[l]{:"},
    literate=
        *{0}{{{\color{numb}0}}}{1}
         {1}{{{\color{numb}1}}}{1}
         {2}{{{\color{numb}2}}}{1}
         {3}{{{\color{numb}3}}}{1}
         {4}{{{\color{numb}4}}}{1}
         {5}{{{\color{numb}5}}}{1}
         {6}{{{\color{numb}6}}}{1}
         {7}{{{\color{numb}7}}}{1}
         {8}{{{\color{numb}8}}}{1}
         {9}{{{\color{numb}9}}}{1}
}

%kopf und fusszeilen
\pagestyle{scrheadings}
\clearscrheadfoot

%kopfzeile
\ihead{\codename}
\ohead{\headmark}
\automark{section}

%fusszeile
\ofoot{\pagemark}
\ifoot{\codename}

\makenoidxglossaries
%
% % Glossareinträge
%

%Kommandos
\newcommand{\codename}{Valaris}


\title{Dokument zur Qualitätssicherung}
\author{Sidney Hansen, Jonas Heinle, Frederik Lingg, Lukas Schölch, Artur Wesner}

\begin{document}
 
	\begin{titlepage}
		\centering
		{\scshape\LARGE Dokument zur Qualitätssicherung\par}
		\vspace{1cm}
		{\scshape\Large Echtzeit Computergrafik in der Spieleentwicklung \par}
		\vspace{1cm}
		{\huge\bfseries \codename \par}
		\vspace{1cm}
		\includegraphics[width=.5\linewidth]{./Bilder/valaris_t_4096.png}
		\par
		{\vspace{1cm}}
		{\Large\itshape Autoren \\}
		{\Large\itshape Sidney Hansen, Jonas Heinle,\\}
		{\Large\itshape Frederik Lingg, Lukas Schölch, Artur Wesner\par}
		
		\vfill
		Projektbetreuer\par
		Alisa Jung \& Tobias Rapp
		
		\vfill
		
		% Bottom of the page
		{\large \today\par}
	\end{titlepage}

	\tableofcontents 
	\pagebreak
	
	
	\pagebreak
	
	\section{Einleitung} 

	Die Aufteilung der Module in InGame-Grafik, InGame-Interface, Assets Generierung und GUI wurde auch in der Qualitätssicherung weiterhin berücksichtigt.
	
	Neben dem Finden und Beheben von Fehlern in der Software war aber auch das Zusammenführen unseres Projekts mit dem Projekt des anderen Teams ein Fokus.
	Dies erwies sich aufgrund von einigen Missverständnissen zum Teil als Herausforderung, welche aber letztendlich bewältigt wurde.

	\section{Qualität}
	\subsection{Definition}
	Qualität ist ein weitreichender Begriff und kann nicht immer gleich verwendet werden.
	Beispielsweise spricht man bei realen Gegenständen wie Textilien von Qualität, wenn man
	die Verarbeitung und Merkmale wie Farbe betrachtet. Bei Lebensmitteln spricht man 
	auch widerrum von einer anderen Qualität. Dementsprechend muss zunächst einmal der 
	Qualitätsbegriff definiert werden, mit dem dieses Software-Projekt bewertet wird. 
	\begin{itemize}
		\item \textbf{Korrektheit}
		\begin{leftbar}[0.9\linewidth]
			Korrektheit ist in Valaris ein ziemlich grober Begriff. Insgesamt müssen 
			alle Aspekte der Streckengenerierung, In-Game Grafik und des Menüs so 
			funktionieren, wie sie in der Entwurfsphase geplant worden sind.
			Folglich müssen alle Tests bezüglich dieser Projektteile möglichst 
			funktionieren, damit hier Qualität gesichert ist. 
		\end{leftbar}
		\item \textbf{Robustheit}
		\begin{leftbar}[0.9\linewidth]
			Unter Robustheit ist zu verstehen, dass Valaris unter verschieden 
			hohen Einstellungen standhaft bleibt. Darunter zählt, dass das Spiel 
			auch unter maximalen Einstellungen angenehm spielbar ist und so dem 
			Benutzer einen hohen Spielspaß liefert. 
		\end{leftbar}
		\item \textbf{Ausfallsicherheit}
		\begin{leftbar}[0.9\linewidth]
			Das Menü sollte auch unter Stressbedingungen, also viele Mouseclicks
			oder Tastatureingaben entsprechend solide reagieren und nicht abstürzen.
			Selbige gilt für die Spielsimulation und deren Grafik. Sollten dennoch 
			Fehler auftreten, so erhält der Benutzer eine entsprechende Meldung.
		\end{leftbar}
		\item \textbf{Benutzerfreundlichkeit}
		\begin{leftbar}[0.9\linewidth]
			Das Spiel sollte intuitiv benutzbar sein. Die Menüführung darf nicht zu 
			kompliziert sein und sollte auch bedienbar sein, wenn das Spiel 
			zum ersten Mal benutzt wird. Falls Fragen bezüglich der Steuerung 
			innerhalb des Spiels auftreten, kann der Benutzer sich über diese informieren,
			wenn er den entsprechenden Punkt im Menü findet.
		\end{leftbar}
		\item \textbf{Effizienz}
		\begin{leftbar}[0.9\linewidth]
			Die Effizienz von Valaris muss daran gemessen werden, wie viel 
			Ressourcen verbraucht werden und wie effizient die Software arbeitet.
			Zu messende Ressourcen sind hier Auslastung des Prozessors, 
			Auslastung der Grafikkarte und der Speicherverbrauch.
		\end{leftbar}
		\item \textbf{Portabilität}
		\begin{leftbar}[0.9\linewidth]
			Hier kommt es auf die Portabilität des Spiels auf andere Systeme an, also
			wie gut die Software sich anpasst auf nutzerspezifische 
			Systemeigenschaften.
		\end{leftbar}
		\item \textbf{Validierbarkeit}
		\begin{leftbar}[0.9\linewidth]
			Validierbarkeit misst die Effektivität der Entwicklung, also wie gut 
			die entstanende Software am Ende zu der anfangs gestellten Aufgabe 
			passt und ob die Umsetzhung zufriedenstellend ist.
		\end{leftbar}
		\item \textbf{Flexibilität}
		\begin{leftbar}[0.9\linewidth]
			Unter Flixibilität ist im Zusammenhang mit Valaris zu verstehen, dass 
			dass der implementierungscode flexibel austauschbar ist. Beispielweise 
			sollen keine \glqq Magic Numbers \grqq{} existieren, was bedeutet, dass 
			Werte in variablen gespeichert werden sollen und die Implementation dann 
			mit diesen arbeitet. 
			Auch soll die Implementierung leicht erweiterbar sein, sobald man 
			weitere Funktionen zum Spiel hinzufügen möchte.
		\end{leftbar}
		\item \textbf{Verständlichkeit}
		\begin{leftbar}[0.9\linewidth]
			Verständlichkeit bezieht sich auf Spielinhalte sowie Code. 
			Bezüglich Spielinhalte kann auf die vorherige Sektion 
			\glqq Benutzerfreundlichkeit \grqq{} verwiesen werden.
			Ein gut verständlicher Code hat JavaDocs zu jeder Funktion und Variable,
			sowie Kommentare innerhalb der Implementation, falls das jeweilige 
			Codestück schwer zu verstehen sein sollte.
		\end{leftbar}
	\end{itemize}
	\subsection{Qualitätsziele}
	Die Qualitätsziele von Valaris liegen in allen Bereichen. 
	Am wenigsten beachtet wurde die Portabilität für andere Systeme, da schon 
	von vornerein festgelegt wurde, dass nicht versucht wird diese zu unterstützen.
	Gemeint ist hier, dass der Fokus eher auf Windows und teilweise auch Linux lag.
	Um die restlichen Ziele zu erreichen, wurden das Werkzeug JUnit benutzt, welches
	automatisierte Tests des Codes mit sich bringt. Außerdem werden so Coverage-
	Tests ebenfalls durchgeführt.
	Um die Effizienz messen zu können, wurde ein eigener Profiler erstellt, welcher
	sich durch die Console anschalten lässt.
	Außerdem wurden bereits in der Pflichtenheftphase manull durchzuführende 
	Testfälle angegeben, um die Qualität und Funktionalität des Spiels zu sichern.
	Diese Testfälle wurden nun durchgegangen und protokolliert.
	Gegen Ende der Qualitätssicherungsphase wurde das Spiel Personen
	vorgestellt, die das Spiel vorher noch nicht gesehen hatten. Diese Testspieler
	spielten ohne Hilfe unserseits, wodurch wichtige Feedbacks gesammelt werden 
	konnten, die diese personen anschließend uns mitgeteilt haben.
	Ein großeß Ziel des Projekt Valaris war unter anderem der Spielspaß. Valaris
	soll dem Spieler am Ende Spaß gemacht haben und um dies zu gewährleisten wurde 
	das Spiel intern viele Male getestet, um so Parameter wie Geschwindigkeit und
	Schwierigkeit zu optimieren.

	\pagebreak

	\section{Änderungen}
	\subsection{InGame - Grafik}

\subsubsection{Neue Features}
\begin{itemize}
    \item \textit{Testfeature}
        \begin{leftbar}[0.9\linewidth]
            Beschreibung des Features
        \end{leftbar}
\end{itemize}

\subsubsection{Änderungen existierender Komponenten}
\begin{itemize}
    \item \textit{Testfeature}
        \begin{leftbar}[0.9\linewidth]
            Beschreibung der Änderung und evtl Begründung
        \end{leftbar}
\end{itemize}

\subsubsection{Nicht implementiert}
\begin{itemize}
    \item \textit{TestNotExistingfeature}
        \begin{leftbar}[0.9\linewidth]
            Alle Methoden die beim Partikeleffekt abhängig von tpf sind wurden nicht
            mit einen JUnit test getestet. Da Ergebnisse abhängig von Framerate macht
            ein klassischer JUnit Test keinen Sinn. Diese wurden allerdings ausführlich
            grafisch auf Funktionalität getestet. 
        \end{leftbar}
\end{itemize}
	\pagebreak
	\subsection{InGame - Interface}

\subsubsection{Neue Features}
\subsubsection{Änderungen existierender Komponenten}
\begin{itemize}
    \item \textit{Tick ist Threadsafe}
        \begin{leftbar}[0.9\linewidth]
            Die Klasse Tick wurde durch das hinzufügen von Synchronisationsmechanismen robuster gemacht.
        \end{leftbar}
\end{itemize}
\subsubsection{Nicht implementiert}
	\subsection{Assets}

\subsubsection{Neue Features}
\begin{itemize}
    \item \textit{Metadaten in AssetDescriptoren}
        \begin{leftbar}[0.9\linewidth]
            Es ist nun möglich Key-Value Paare im AssetDescriptor zu definieren und diese
            zur Laufzeit durch das AssetControl aus zu lesen
        \end{leftbar}
\end{itemize}

\subsubsection{Änderungen existierender Komponenten}
\begin{itemize}
    \item \textit{Laden von Controls}
        \begin{leftbar}[0.9\linewidth]
            Das AssetControl kann nun ausschließlich Controls vom Typ AbstractAssetControl laden,
            ist dadurch aber in der Lage geladene Controls über das zerstören ihres Parents zu benachrichtigen.
        \end{leftbar}
\end{itemize}
\subsubsection{Nicht implementiert}
	\subsection{Generierung}

\subsubsection{Neue Features}
\begin{itemize}
    \item \textit{Bepflanzung der Kuppeln}
        \begin{leftbar}[0.9\linewidth]
            Es können nun für ein bestimmtes Biom eine Menge von Bäumen(bzw. ähnliche Pflanzen) definiert werden,
            welche dann in einer entsprechenden Kuppel platziert werden. Bei der Platzierung werden die Objekte
            anhand einer, im \textit{DomeAssetGenerator} definierten, Dichtefunktion verschoben, wie stark 
            kann durch einen Parameter pro Biom bestimmt werden.
            Desweiteren kann man pro Biom unter anderem die Asset-Skalierung und Farben der Blätter festlegen.
        \end{leftbar}

    \item \textit{Assets als Kuppeldach}
        \begin{leftbar}[0.9\linewidth]
            Es ist nun möglich Assets als Dach für eine Kuppel zu verwenden. Es können auch verschiedene Assets
            definiert werden, aus denen dann ein zufälliges(abhängig vom Seed) ausgesucht wird.
        \end{leftbar}
\end{itemize}

\subsubsection{Änderungen existierender Komponenten}
\begin{itemize}
    \item \textit{Methode \textit{generateScenegraph()} parallelisiert}
        \begin{leftbar}[0.9\linewidth]
            \textit{GenerateScenegraph()} der Klasse Map, ruft nun die \textit{generateScenegraph()} Methoden
            der ISceneItem parallelisiert auf und beschleunigt so den Generierungsprozess deutlich.
        \end{leftbar}
    
    \item \textit{'Colring' von Landschaften verschönert}
        \begin{leftbar}[0.9\linewidth]
            Es entstehen nun keine ungewollten (bunte) Dreiecksstrukturen mehr, da das festlegen der 'Oberflächen'
            an einem Punkt optimiert wurde.
        \end{leftbar}

    \item \textit{Verbesserte Räume}
        \begin{leftbar}[0.9\linewidth]
            Ein weiterer Raumtyp wurde hinzugefügt, sowie die Asset-Platzierung in räumen optimiert.
        \end{leftbar}

    \item \textit{Unterschiedliche Strecken}
        \begin{leftbar}[0.9\linewidth]
            Unterschiedliches Straßen Assets in Tunneln, Kuppeln und Räumen.
        \end{leftbar}

    \item \textit{Neue Biome}
        \begin{leftbar}[0.9\linewidth]
            Es wurden 4 fertige Biome definiert, die sich in den Punkten Untergrund, Farben und platzierten Objekten
            unterscheiden.
        \end{leftbar}

    \item \textit{Kuppeldach wurde durch Asset ersetzt}
        \begin{leftbar}[0.9\linewidth]
            Das Dach einer Kuppel besteht nun aus einem, in Blender erstelltem, Asset.
        \end{leftbar}

    \item \textit{Definition von 3 Kartentypen}
        \begin{leftbar}[0.9\linewidth]
            Es wurden ein kleines, ein mittleres und ein großes Biom definiert, die sich in der Anzahl der Kuppeln, 
            Tunneln und Räumen unterscheiden. Ebenso sind di möglichen Größen eines solchen \textit{ISceneItems} unterschiedlich.
        \end{leftbar}
\end{itemize} 
	\subsection{GUI}

\subsubsection{IConfig}
\begin{itemize}
    \item \textit{Umbenannt zu ConfigManager und kein Interface mehr}
        \begin{leftbar}[0.9\linewidth]
            Hier lag ein Entwurfsfehler vor. Diese Klasse stellt nun Funktionen zum Lesen und 
            Schreiben einer Datei bereit. Sie enthält für jede Methode einen Pfad, wo sich 
            die Datei befinden sollte.
        \end{leftbar}
\end{itemize}

\subsubsection{InternalGameConfig, SeedConfig und GraphicsConfig}
\begin{itemize}
    \item \textit{Implementiert nicht mehr IConfig}
        \begin{leftbar}[0.9\linewidth]
            Da die Klasse statische Methoden enthält, implementiert sie nicht mehr das 
            Interface IConfig.
        \end{leftbar}
    \item \textit{Methode getConfig() hinzugefügt}
        \begin{leftbar}[0.9\linewidth]
            Gibt die Konfiguration zurück und ruft bei die getConfig()
            Methode vom ConfigManager auf.
        \end{leftbar}
    \item \textit{Methode saveConfig() hinzugefügt}
        \begin{leftbar}[0.9\linewidth]
            Speichert die Konfiguration und ruft bei die 
            saveConfig() Methode vom ConfigManager auf.
        \end{leftbar}
    \item \textit{Methode reset() hinzugefügt}
        \begin{leftbar}[0.9\linewidth]
            Setzt die Config zurück durch Überschreiben mit der Default-Datei.
        \end{leftbar}
\end{itemize}

\subsubsection{OwnScreenController}
\begin{itemize}
    \item \textit{Klasse hinzugefügt}
        \begin{leftbar}[0.9\linewidth]
            OwnScreenController implementiert von ScreenController und wird von jedem Menü-
            Bildschirm geerbt. Es enthält Funktionen und Attribute, die jeder Bildschirm 
            braucht und reduziert somit Code-Duplizierung. Jedem OwnScreenController wird 
            außerdem eine Referenz auf den NiftyAppState und die anderen OwnScreenController
            mitgegeben, so dass diese leicht aufrufbar sind.
        \end{leftbar}
\end{itemize}

\subsubsection{NiftyAppState}
\begin{itemize}
    \item \textit{Gesamtes Entwurfskonzept der Klasse geändert}
        \begin{leftbar}[0.9\linewidth]
            Ein NiftyAppState ist nun eine AppState, welcher alle die ihm mitgegebenen 
            OwnScreenController in einer Nifty Instanz initialisiert und speichert.
            Die OwnScreenController innerhalb eines NiftyAppStates haben Zugriff auf 
            einander und können so leicht mit nifty.goToScreen() aufgerufen werden.
            Ist dies nicht möglich, so erstellt ein OwnScreenController ein neues 
            NiftyAppState-Objekt mit den jeweiligen OwnScreenControllern und fügt dieses
            an den AppStateManager an.
            Wird der NiftyAppState vom AppStateManager entfernt, so entfernt dieser alle 
            initialisierten Sachen wie z.B den Processor für das NiftyDisplay.
        \end{leftbar}
\end{itemize}

\subsubsection{SplashScreen}
\begin{itemize}
    \item \textit{Klasse wurde hinzugefügt}
        \begin{leftbar}[0.9\linewidth]
            Zeigt einen SplashScreen bei Start des Programms an.
        \end{leftbar}
\end{itemize}

\subsubsection{MainmenuScreen}
\begin{itemize}
    \item \textit{Implementiert nicht mehr ScreenController, erbt dafür von OwnScreenController}
        \begin{leftbar}[0.9\linewidth]
            Es wurde eine eigene ScreenController Klasse hinzugefügt.
        \end{leftbar}
    \item \textit{Methode openMultiplayerMode() entfernt}
        \begin{leftbar}[0.9\linewidth]
            Wurde nicht mehr benötigt, da die Mehrspieler-Funktion inklusive des 
            Lobby Bildschirmes entfernt wurde.
        \end{leftbar}
\end{itemize}

\subsubsection{AbstractEditGameSettings}
\begin{itemize}
    \item \textit{Klasse entfernt}
        \begin{leftbar}[0.9\linewidth]
            Diese Klasse existierte ursprünglich, da der Lobby Bildschirm und der 
            Einzelspielerbildschirm viele gleiche Methoden brauchten. Da der Lobby 
            Bildschirm entfernt wurde, wurden alle Funktionen hierraus in die Klasse 
            SinglePlayerModeScreen verschoben.
        \end{leftbar}
\end{itemize}

\subsubsection{LobbyScreen}
\begin{itemize}
    \item \textit{Klasse entfernt}
        \begin{leftbar}[0.9\linewidth]
            Wurde nicht mehr benötigt, da die Mehrspieler-Funktion inklusive des 
            Lobby Bildschirmes entfernt wurde.
        \end{leftbar}
\end{itemize}

\subsubsection{SinglePlayerModeScreen}
\begin{itemize}
    \item \textit{Implementiert nicht mehr ScreenController, erbt dafür von OwnScreenController}
        \begin{leftbar}[0.9\linewidth]
            Es wurde eine eigene ScreenController Klasse hinzugefügt.
        \end{leftbar}
    \pagebreak
    \item \textit{Methoden aus AbstractEditGameSettings hier eingefügt}
        \begin{leftbar}[0.9\linewidth]
            AbstractEditGameSettings wurde entfernt und dementsprechend sind dessen Methoden 
            hier eingesetzt worden.
        \end{leftbar}
    \item \textit{Methode editCurrentSeed(String, DropDownSelectionChangedEvent) hinzugefügt}
        \begin{leftbar}[0.9\linewidth]
            Wird aufgerufen sobald eine Auswahl im DropDown-Menü getroffen wird.
        \end{leftbar}
    \item \textit{Methode editCurrentSeedNumber(String, TextFieldChangedEvent) hinzugefügt}
        \begin{leftbar}[0.9\linewidth]
            Wird aufgerufen sobald die Seed Nummer bearbeitet wird.
        \end{leftbar}
    \item \textit{Methode editKICount(String, TextFieldChangedEvent) hinzugefügt}
        \begin{leftbar}[0.9\linewidth]
            Wird aufgerufen sobald die KI Anzahl bearbeitet wird.
        \end{leftbar}
    \item \textit{Methode saveSeedName() hinzugefügt}
        \begin{leftbar}[0.9\linewidth]
            Wird aufgerufen sobald der Seed Name mit Enter bestätigt, oder der Speicher-Button
            rechts daneben geklickt wird.
        \end{leftbar}
    \item \textit{Methode editSeednameTextfieldVisibility() hinzugefügt}
        \begin{leftbar}[0.9\linewidth]
            Wird aufgerufen sobald der Bearbeiten-Button neben dem DropDown-Menü gedrückt wird.
        \end{leftbar}
    \item \textit{Methode saveSeedNumber() hinzugefügt}
        \begin{leftbar}[0.9\linewidth]
            Wird aufgerufen sobald die Seed Nummer mit Enter bestätigt, oder der Speicher-Button
            rechts daneben geklickt wird.
        \end{leftbar}
\end{itemize}

\subsubsection{SettingsmenuScreen}
\begin{itemize}
    \item \textit{Implementiert nicht mehr ScreenController, erbt dafür von OwnScreenController}
        \begin{leftbar}[0.9\linewidth]
            Es wurde eine eigene ScreenController Klasse hinzugefügt.
        \end{leftbar}
    \item \textit{Methode editCurrentResolution(String, DropDownSelectionChangedEvent) hinzugefügt}
        \begin{leftbar}[0.9\linewidth]
            Wird aufgerufen sobald eine Auswahl im Auflösungs-DropDown getroffen wird.
        \end{leftbar}
    \item \textit{Methode editCurrentLanguage(String, DropDownSelectionChangedEvent) hinzugefügt}
        \begin{leftbar}[0.9\linewidth]
            Wird aufgerufen sobald eine Auswahl im Sprache-DropDown getroffen wird.
        \end{leftbar}
    \item \textit{Methode editCurrentKIDifficulty(String, DropDownSelectionChangedEvent) hinzugefügt}
        \begin{leftbar}[0.9\linewidth]
            Wird aufgerufen sobald eine Auswahl im KI-Schwiriegskeits-DropDown getroffen wird.
        \end{leftbar}
    \item \textit{Methode editCurrentQuality(String, DropDownSelectionChangedEvent) hinzugefügt}
        \begin{leftbar}[0.9\linewidth]
            Wird aufgerufen sobald eine Auswahl im Qualitäts-DropDown getroffen wird.
        \end{leftbar}
    \item \textit{Methode toggleFullscreen(String, CheckBoxStateChangedEvent) hinzugefügt}
        \begin{leftbar}[0.9\linewidth]
            Wird aufgerufen sobald eine Auswahl in der Vollbild-Checkbox getroffen wird.
        \end{leftbar}
    \item \textit{Methode showGraphicsTab() hinzugefügt}
        \begin{leftbar}[0.9\linewidth]
            Zeigt den Grafik Einstellungs-Tab an.
        \end{leftbar}
    \item \textit{Methode showInternTab() hinzugefügt}
        \begin{leftbar}[0.9\linewidth]
            Zeigt den Internen Einstellungs-Tab an.
        \end{leftbar}
    \item \textit{Methode resetSettings() hinzugefügt}
        \begin{leftbar}[0.9\linewidth]
            Setzt die Einstellungen zurück.
        \end{leftbar}
\end{itemize}

\subsubsection{LoadingScreen}
\begin{itemize}
    \item \textit{Implementiert nicht mehr ScreenController, erbt dafür von OwnScreenController}
        \begin{leftbar}[0.9\linewidth]
            Es wurde eine eigene ScreenController Klasse hinzugefügt.
        \end{leftbar}
    \item \textit{Folgende Methoden entfernt}
        \begin{itemize}
            \item openHUD() : HUD wird an anderer Stelle geöffnet.
            \item generateMap() : Funktionalität wude in startGame() verschoben.
            \item getMapBody() : Nicht mehr benötigt.
        \end{itemize}
\end{itemize}

\subsubsection{LoadingUtil}
\begin{itemize}
    \item \textit{Klasse hinzugefügt}
        \begin{leftbar}[0.9\linewidth]
            In dieser Util Klasse existieren Funktionen, die vom LoadingScreen aufgerufen 
            werden, um alle benötigten Sachen zu laden. Mittels Callbacks wird der 
            ladebalken entsprechend gesetzt.
        \end{leftbar}
\end{itemize}

\subsubsection{Hud}
\begin{itemize}
    \item \textit{Implementiert nicht mehr ScreenController, erbt dafür von OwnScreenController}
        \begin{leftbar}[0.9\linewidth]
            Es wurde eine eigene ScreenController Klasse hinzugefügt.
        \end{leftbar}
    \item \textit{Alle Funktionen entfernt/ersetzt}
        \begin{leftbar}[0.9\linewidth]
            Die Minimap wurde nicht im Hud integriert und die restlichen Funktionen wurden 
            ersetzt bzw entfernt. Funktionen um auf andere Bildschirme du wechseln wurde an 
            eine andere Stelle verlagert.
        \end{leftbar}
    \pagebreak
    \item \textit{Methode updateRank(int) hinzugefügt}
        \begin{leftbar}[0.9\linewidth]
            Aktualisiert den angezeigten Rang im Hud.
        \end{leftbar}
    \item \textit{Methode updateTime(float) hinzugefügt}
        \begin{leftbar}[0.9\linewidth]
            Aktualisiert die angezeigte Zeit im Hud.
        \end{leftbar}
    \item \textit{Methode updateCountdown(String) hinzugefügt}
        \begin{leftbar}[0.9\linewidth]
            Aktualisiert den angezeigten Countdown im Hud.
        \end{leftbar}
    \item \textit{Methode updateItem(String) hinzugefügt}
        \begin{leftbar}[0.9\linewidth]
            Aktualisiert das angezeigte Item im Hud.
        \end{leftbar}
\end{itemize}

\subsubsection{ErrorScreen}
\begin{itemize}
    \item \textit{Klasse hinzugefügt}
        \begin{leftbar}[0.9\linewidth]
            Diese Klasse dient dazu, eventuell auftretende Exceptions/Fehler während dem Spiel 
            anzeigen zu können mittels eines Popups.
        \end{leftbar}
\end{itemize}

\subsubsection{PausemenuScreen}
\begin{itemize}
    \item \textit{Implementiert nicht mehr ScreenController, erbt dafür von OwnScreenController}
        \begin{leftbar}[0.9\linewidth]
            Es wurde eine eigene ScreenController Klasse hinzugefügt.
        \end{leftbar}
\end{itemize}

\subsubsection{EndmenuScreen}
\begin{itemize}
    \item \textit{Implementiert nicht mehr ScreenController, erbt dafür von OwnScreenController}
        \begin{leftbar}[0.9\linewidth]
            Es wurde eine eigene ScreenController Klasse hinzugefügt.
        \end{leftbar}
    \item \textit{Methode updatePlayerEndTime() entfernt}
        \begin{leftbar}[0.9\linewidth]
            Ersetzt durch folgende Methode.
        \end{leftbar}
    \item \textit{Methode updatePlayerEndTimeAndColor(int, float, String) hinzugefügt}
        \begin{leftbar}[0.9\linewidth]
            Aktualisiert die Zeit und Farbe des Spielers.
        \end{leftbar}
\end{itemize}

	\pagebreak
	\section{Fehlerbehebungen}
	\subsection{InGame - Grafik}
\subsubsection{Behoben}

\begin{itemize}
    \item wrongColor
    \includegraphics[width=0.7\textwidth]{Grafik/pictures/wrongColor.png} %70% der Textbreite
    \item wrongRampFac
    \includegraphics[width=0.7\textwidth]{Grafik/pictures/rampFac.png} %70% der Textbreite
    \item Achsendreher
    \includegraphics[width=0.7\textwidth]{Grafik/pictures/Achsendreher.png} %70% der Textbreite
\end{itemize}

\subsubsection{Nicht behoben}
\begin{itemize}
    \item Nothing
    
\end{itemize}
	\subsection{InGame - Interface}
\subsubsection{Behoben}
\begin{itemize}
    \item Tick-Struckturen waren anfangs nicht Threadsafe, dies wurde durch Einfügen von Locks behoben.
\end{itemize}

\subsubsection{Nicht behoben}
\begin{itemize}
    \item Selten gehen vermutlich bei synchronisationsstrukturen TickEvents verloren, wodurch ein unzulässiger Zustand im TickProcessor hervorgerufen wird,
    was zum Abbruch des Rennens führt.
\end{itemize}
	\subsection{Assets}
\subsubsection{Behoben}
\begin{itemize}
    \item Materialien wurden gesetzt nachdem Controls geladen wurden
\end{itemize}
\subsubsection{Nicht behoben}
	\pagebreak
	\subsection{Generierung}
\subsubsection{Behoben}
\begin{itemize}
    \item Kuppelwand-Teilstücke manchmal falsch generiert.
    \item 'Löcher' unter Ein-/Ausgängen der Kuppelwand.
    \item Coloring funktionierte nicht wie erwartet.
    \item Letze Reihe an GridVertices wurde (selten) nicht der Straße angepasst.
    \item Start- und Endraum wurde nicht erstellt.
    \item Straßen teilweise falsch skaliert.
    \item Straße in Kuppeln und Tunneln war 'kantig'.
    \item Offset zwsichen Kuppelausgang und Tunneleingang.
    \item Offset zwsichen Raumeingang und Tunnelausgang bzw. Raumausgang und Tunneleingang.
    \item Zu große Skalierung der Assets in Räumen.
    \item Jedes zweite Licht wurde entfernt.
    \item Dimensionen des ersten RoadCursurs der Kuppel wurden nicht richtig gesetzt (Türdimensionen passten nicht).
    \item Straßengenerierung der Kuppel ging von einer kugelförmigen Wand aus.
    \item Diverse Umrechnungsprobleme beim MapBody.
    \item Startpositionen wurden nicht übergeben.
\end{itemize}

\subsubsection{Nicht behoben}
\begin{itemize}
    \item Straße in Räumen ist teilweise 'unförmig'.
    \item Bäume stehen nicht optimal auf dem Boden (aufgrund der Triangulation).
\end{itemize}
	\pagebreak
	\subsection{GUI}
\subsubsection{Behoben}
\begin{itemize}
    \item Konfigurationspfade wurden erneuert. Sollte diese trotzdem nicht gefunden
    werden, so wird eine Default-Konfiguration geladen.
    \item Ändert der Benutzer manuell die Konfigurationsdatei und etwas entspricht
    keinem korrekten Wert, so wird die Konfiguration zurückgesetzt.
    \item Die Schrift des aktuellen Stands im Ladebalken war manchmal verschoben.
    \item Es wird nun die korrekte Anzahl an Fahrzeugen an das Physik-Team 
    übergeben.
    \item Der Countdown wird nun korrekt angehalten, wenn das Spiel pausiert.
    \item Der Wechsel zwischen Pause und Spiel erfolgt nun in jedem Zustand des 
    Spieles erfolgreich. Eingeblendete Informationen, wie die Zeit, werden 
    erfolgreich pausiert.
    \item Nifty Objekte reagierten immer noch auf Events, obwohl dessen 
    zugehöriger NiftyAppState vom StateManager entfernt wurden.
    \item Seeds werden nun unter dem Format \glqq SeedZahl\#Streckentyp \grqq{} gespeichert.
    \item Die Farbe eines Spielers wird nun korrekt im End-Menü angezeigt.
    \item Änderte man den Namen eines Seed's, wurde der alte Seed beibehalten 
    und ein neuer Seed mit dem neuen Namen angelegt.
    \item Die Auflösung konnte nicht während des Vollbildmodus geändert werden.
    \item Beim Sprachwechsel wurde die aktuelle Drop-Down Auswahl nicht korrekt 
    gesetzt, da die Inhalte sich verändert haben. Dies ist nun für jedes 
    Drop-Down behoben. 
\end{itemize}

\subsubsection{Nicht behoben}
\begin{itemize}
    \item Die Sprache kann während dem Zurücksetzen der Konfiguration nicht mit 
    geändert werden, daher wird eine entsprechende Meldung angezeigt.
\end{itemize}
	\pagebreak

	

	\section{Neu hinzugekommene Tests}

	\subsection{InGame - Grafik}

\subsubsection{Automatische Tests}
\paragraph{TestKlasse1}
\begin{itemize}
    \item \textit{Testtestfall}
        \begin{leftbar}[0.9\linewidth]
			Partikeleffekt wird hinzugefügt zum Spiel.
			Partikeleffekt wird angezeigt.
        \end{leftbar}
\end{itemize}

\subsubsection{Manuelle Tests}
	\begin{itemize} [label={}]
		\item \textbf{ET0XX}\par 
			\textbf{STAND} Fahrzeug fährt\\
			\textbf{AKTION} Explosion als DgO\\
            \textbf{RESULTAT} Explosion ausgeführt.\\\par
            
       

	\end{itemize}

	\input{Interface/Interface_NewTests}
	\subsection{Assets}
\subsubsection{Automatische Tests}
\subsubsection{Manuelle Tests}

	\subsection{Generierung}

\subsubsection{Automatische Tests}
\paragraph{BiomMeadowColoringTest}
\begin{itemize}
    \item \textit{makeColorRoad}
        \begin{leftbar}[0.9\linewidth]
			Testet die Färbung zweier Dreiecke, bei denen ein Fahrbahnuntergrund zu sehen sein soll.
			Jeder Farbkanal enthält den gleichen Farbwert: 0.1f.
		\end{leftbar}
		
	\item \textit{makeColorRock}
        \begin{leftbar}[0.9\linewidth]
			Testet die Färbung zweier Dreiecke, bei denen ein Stein zu sehen sein soll.
			Jeder Farbkanal enthält den gleichen Wert: 0f.
		\end{leftbar}
		
	\item \textit{makeColorSnow}
        \begin{leftbar}[0.9\linewidth]
            Testet die Färbung zweier Dreiecke, bei denen Schnee zu sehen sein soll.
			Jeder Farbkanal enthält den gleichen Wert: 1f.
		\end{leftbar}
		
	\item \textit{makeColorMeadow}
		\begin{leftbar}[0.9\linewidth]
			Testet die Färbung zweier Dreiecke, bei denen Wiese zu sehen sein soll.
			Jeder Farbkanal enthält unterschiedliche Werte zwischen 0f und 1f.
		\end{leftbar}

	\item \textit{makeColorSnowOverMeadow}
        \begin{leftbar}[0.9\linewidth]
			Testet die Färbung zweier Dreiecke, bei denen nicht alle Eckpunkte die gleichen Eigenschaftenm aufweisen,
			einige Eckpunkte 'bestehen aus' Wiese andere aus Schnee.
			Die Farbkanäle enthalten die oben genannten Werte, je für Schnee bzw. Wiese.
		\end{leftbar}
		
	\item \textit{makeColorMeadowOverRock}
        \begin{leftbar}[0.9\linewidth]
			Testet die Färbung zweier Dreiecke, bei denen nicht alle Eckpunkte die gleichen Eigenschaftenm aufweisen,
			einige Eckpunkte 'bestehen aus' Wiese andere aus Stein.
			Die Farbkanäle enthalten die oben genannten Werte, je für Stein bzw. Wiese.
        \end{leftbar}
\end{itemize}

\paragraph{MapBodyTest}
\begin{itemize}
    \item \textit{test0Calc3DTransform}
        \begin{leftbar}[0.9\linewidth]
			Testet die Transformat von Positionen aus dem RoadModel
			in Weltkoordinaten.
		\end{leftbar}
	\item \textit{test1Calc3DTransform}
        \begin{leftbar}[0.9\linewidth]
            Testet die Transformat von Positionen aus dem RoadModel
			in Weltkoordinaten.
		\end{leftbar}
	\item \textit{test2Calc3DTransform}
        \begin{leftbar}[0.9\linewidth]
            Testet die Transformat von Positionen aus dem RoadModel
			in Weltkoordinaten.
		\end{leftbar}
\end{itemize}

\paragraph{RoadGeneratorTest}
\begin{itemize}
	\item \textit{generateSphereRoad}
        \begin{leftbar}[0.9\linewidth]
			Testet ob die Generierung einer Strecke für eine Kuppel erfolgreich war.
			Testet die Koorektheit der Rotation der RoadCursor Instanzen.
			Testet ob die Minimaldistanz zwischen den RoadCursor Instanzen eingehalten.
		\end{leftbar}
	\item \textit{generateConnectingRoad}
        \begin{leftbar}[0.9\linewidth]
            Testet ob die Generierung einer Strecke für einen Tunnel erfolgreich war.
			Testet die Koorektheit der Rotation der RoadCursor Instanzen.
			Testet ob die Minimaldistanz zwischen den RoadCursor Instanzen eingehalten.
		\end{leftbar}
	\item \textit{generateCuboidRoad}
        \begin{leftbar}[0.9\linewidth]
            Testet ob die Generierung einer Strecke für einen Raum erfolgreich war.
			Testet die Koorektheit der Rotation der RoadCursor Instanzen.
			Testet ob die Minimaldistanz zwischen den RoadCursor Instanzen eingehalten.
		\end{leftbar}
\end{itemize}

\paragraph{DynamicRoomGeneratorTest}
\begin{itemize}
	\item \textit{test0GenerateRoom}
        \begin{leftbar}[0.9\linewidth]
            Testet ob die Raumgenerierung unter bestimmten Parametern erfolgreich ist.
		\end{leftbar}
	\item \textit{test1GenerateRoom}
        \begin{leftbar}[0.9\linewidth]
            Testet ob die Raumgenerierung unter bestimmten Parametern erfolgreich ist.
		\end{leftbar}
	\item \textit{test2GenerateRoom}
        \begin{leftbar}[0.9\linewidth]
            Testet ob die Raumgenerierung unter bestimmten Parametern erfolgreich ist.
		\end{leftbar}
	\item \textit{testGenerateStartRoom}
        \begin{leftbar}[0.9\linewidth]
			Testet ob die Raumgenerierung unter bestimmten Parametern erfolgreich ist.
			Testet ob die Richtige Anzahl an Startpositionen Platziert wurde.
		\end{leftbar}
	\item \textit{testGenerateEndRoom}
        \begin{leftbar}[0.9\linewidth]
			Testet ob die Raumgenerierung unter bestimmten Parametern erfolgreich ist.
			Testet ob die Ziellinie korrekt platziert wurde.
		\end{leftbar}
\end{itemize}


\paragraph{TunnelGeneratorTest}
\begin{itemize}
	\item \textit{testGenerate}
        \begin{leftbar}[0.9\linewidth]
			Testet ob die TunnelGenerierung mit unterschiedlichen Parametern erfolgreich ist.
			Testet ob der letzte RoadCursor in Position und Rotation nur minimal vom "Ziel"-RoadCursor abweicht.
		\end{leftbar}
\end{itemize}


\paragraph{RoadTest}
\begin{itemize}
	\item \textit{testUpdateDirections}
        \begin{leftbar}[0.9\linewidth]
			Testet ob die Methode testUpdateDirections(...) die Rotation der RoadCursor korrekt korrigiert.
		\end{leftbar}
	\item \textit{attach}
        \begin{leftbar}[0.9\linewidth]
			Testet ob die Methode attach(...) die Positionen der RoadCursor Instanzen der
			angehängten Road korrekt transformiert.
		\end{leftbar}
\end{itemize}

\paragraph{RoadCursorTest}
\begin{itemize}
	\item \textit{testPosition}
        \begin{leftbar}[0.9\linewidth]
			Testet die Position der RoadCursor:
			Nach dem Erstellen durch unterschiedliche Konstruktoren,
			Nach Anwendung von Rotation durch die Methode applyRotation(...),
			Nach Anwendung von Transformation durch die Methode applyTransformation(...);
		\end{leftbar}
	\item \textit{testDirection}
        \begin{leftbar}[0.9\linewidth]
			Testet den Richtungsvektor der RoadCursor:
			Nach dem Erstellen durch unterschiedliche Konstruktoren,
			Nach Anwendung von Rotation durch die Methode applyRotation(...),
			Nach Anwendung von Transformation durch die Methode applyTransformation(...);
		\end{leftbar}
	\item \textit{testRight}
        \begin{leftbar}[0.9\linewidth]
			Testet den Rechtsvektor der RoadCursor:
			Nach dem Erstellen durch unterschiedliche Konstruktoren,
			Nach Anwendung von Rotation durch die Methode applyRotation(...),
			Nach Anwendung von Transformation durch die Methode applyTransformation(...);
		\end{leftbar}
	\item \textit{testNormal}
        \begin{leftbar}[0.9\linewidth]
			Testet den Normalenvektor der RoadCursor:
			Nach dem Erstellen durch unterschiedliche Konstruktoren,
			Nach Anwendung von Rotation durch die Methode applyRotation(...),
			Nach Anwendung von Transformation durch die Methode applyTransformation(...);
		\end{leftbar}
	\item \textit{testXZAngle}
        \begin{leftbar}[0.9\linewidth]
			Testet die Berechnung des Winkels in der Ebene der RoadCursor.
		\end{leftbar}
	\item \textit{testYAngle}
        \begin{leftbar}[0.9\linewidth]
			Testet die Berechnung des Steigungs Winkels der RoadCursor.
		\end{leftbar}
	\item \textit{testTiltAngle}
        \begin{leftbar}[0.9\linewidth]
			Testet ob der Kippwinkel des RoadCursors 0 ist.
		\end{leftbar}
	\item \textit{testRoadCursorProperty}
        \begin{leftbar}[0.9\linewidth]
			Testet ob der RoadCursor die ihm zugewiesene RoadCursorProperty hat.
		\end{leftbar}
	\item \textit{testRoadFrame}
        \begin{leftbar}[0.9\linewidth]
			Testet ob die Methode testRoadFrame(...) korrekt funktioniert.
		\end{leftbar}
\end{itemize}

\subsubsection{Manuelle Tests}

	Zusätzlich zu den oben genannten JUnit-Tests wurde zum größten Teil manuell getestet, da optische Eindrücke 
	nicht automatisiert getestet werden können. Hierzu wurde für jeden betrachteten Teil ggf. die 
	\textit{GenerationConfig} kurzzeitig angepasst und dann das Spiel mit der DummySimulation ausgeführt, um so das
	entsprechende Element auf Sicht zu prüfen. Hierzu wurden die betrachteten Teile vorallem möglichst klein gehalten
	und manchmal auch beispielsweise Screenshots als Hilfestellung verwendet. Unten aufgelistet sind beispielhafte 
	Testabläufe, welche noch nicht im Pflichtenheft definiert wurden.

	\begin{itemize} [label={}]
		\item \textbf{ET120}\par %Bitte im Pflichtenheft nachschauen nach der nächsten Nummer
			\textbf{STAND} Benutzer im Hauptmenü.\\
			\textbf{AKTION} Generierung mit beleuchteten Räumen.\\
            \textbf{RESULTAT} Beleuchtung wirkt natürlich.\\\par
            
        \item \textbf{ET130}\par %Bitte im Pflichtenheft nachschauen nach der nächsten Nummer
			\textbf{STAND} Generierung, mit Kuppel und darin leigender Straße, abgeschlossen.\\
			\textbf{AKTION} Betrachtung der auf die Straße angepassten Landschaft.\\
			\textbf{RESULTAT} Landschaft unterhalb der Straße wird, in allen Punkten, 
			genau unter die Straße gesetz.\\\par

	\end{itemize}
	\pagebreak
	\subsection{GUI}

\subsubsection{Automatische Tests}
\begin{itemize}
    \item \textit{Alle Klassen}
        \begin{leftbar}[0.9\linewidth]
			Die Testmethoden des Menüs wurde lediglich an geforderten Stellen 
			erneuert, sofern sich Code geändert hat.
        \end{leftbar}
\end{itemize}

\subsubsection{Manuelle Tests}
	\begin{itemize} [label={}]
		\item \textbf{ET203}\par 
			\textbf{STAND} Startmenü\\
			\textbf{AKTION} Spiel starten nach auswählen des Streckentyps\\
			\textbf{RESULTAT} Strecke hat folgende Anzahl an Kuppeln, je nach Streckentyp:
			Klein: 1-2. Normal: 3-5. Groß: 5-7.\\\par
            
        \item \textbf{ET241}\par 
			\textbf{STAND} Laufendes Rennen\\
			\textbf{AKTION} Im Ziel ankommen, warten bis sich das Ziel-Menü öffnet
			und anschließend Übeprüfung der Inhalte.\\
			\textbf{RESULTAT} Die Zeiten und Platzierungen der Fahrer sollten übereinstimmen. 
			Die Farben sollten nur grob übereinstimmen, da die Anzeige der Farbe 
			des Fahrzeugs während dem Spiel auf Grund von einigen Parametern wie z.B. Licht, 
			für den Benutzer anders erscheinen werden als im Menü.\\\par

	\end{itemize}

	\pagebreak

	\section{Testergebnisse}
	\subsection{InGame - Grafik}

\begin{tabularx}{0.9\textwidth}{p{0.25\textwidth}|p{0.15\textwidth}||p{0.25\textwidth}|p{0.15\textwidth}}
    \textbf{Testnummer} & \textbf{Funktional} & \textbf{Testnummer} & \textbf{Funktional}\\
    \hline
    ET010	& Ja& ET011	& Nein	\\
    \hline              
    ET012   & Ja   & ET020   & Ja  \\      
    \hline            
    ET030   &Ja & Et040   &Ja \\
\end{tabularx}
	\subsection{InGame - Interface}

\begin{tabularx}{0.9\textwidth}{p{0.25\textwidth}|p{0.15\textwidth}||p{0.25\textwidth}|p{0.15\textwidth}}
    \textbf{Testnummer} & \textbf{Funktional} & \textbf{Testnummer} & \textbf{Funktional}\\
    \hline
    ET010	& Ja		& ET012		& Ja\\
    \hline
    ET011	& Ja
\end{tabularx}
	\subsection{Assets}
Im Pflichtenheft wurden keine Tests für dieses Modul definiert,
da die Notwendigkeit des Moduls erst später festgestellt wurde.
	\subsection{Generierung}

\begin{tabularx}{0.9\textwidth}{p{0.25\textwidth}|p{0.15\textwidth}||p{0.25\textwidth}|p{0.15\textwidth}}
    \textbf{Testnummer} & \textbf{Funktional} & \textbf{Testnummer} & \textbf{Funktional}\\
    \hline
    ET100	& Ja		& ET111		& Ja\\
    \hline
    ET101	& Ja		& ET112	    & Ja\\
    \hline
    ET102	& Ja		& ET112		& Ja\\
    \hline
    ET110	& Ja		& ET114	    & Ja\\
\end{tabularx}
	\subsection{GUI}

\begin{tabularx}{0.9\textwidth}{p{0.25\textwidth}|p{0.15\textwidth}||p{0.25\textwidth}|p{0.15\textwidth}}
    \textbf{Testnummer} & \textbf{Funktional} & \textbf{Testnummer} & \textbf{Funktional}\\
    \hline
    ET200	& Ja		& ET220		& Ja\\
    \hline
    ET201	& Ja		& ET230		& Ja\\
    \hline
    ET202	& Entfernt		& ET240		& Ja\\
    \hline
    ET203	& Ja		& ET241		& Ja\\
    \hline
    ET210	& Ja		& 		& \\
\end{tabularx}
	\pagebreak

	\subsection{UserTests}
	Im Zuge der Qualitätsverbesserung wurde eine spielbare Version an Personen 
	außerhalb des Projektes versendet. Ziel war, dadurch Erfahrungen und 
	Feedback zu sammeln aus der Sicht einer Person, die nicht an der 
	Implementierung beteiligt war.
	Diesen Personen wurde als einzige Aufgabe gestellt, das Spiel zu spielen.
	Sollten Fehler auftreten oder sonstige Schwierigkeiten, wo wurde 
	um eine Mitteilung gebeten. Sobald die personen dann fertig waren, wurde 
	um ein abschließendes Feedback gebeten, welche Punkte positiv und negativ 
	waren.
	Viele der Testpersonen hatten angemerkt, dass die Strecke etwas dunkel sei, 
	vorallem in den Tunneln und am Anfang der Strecke. Dies wurde von uns dann 
	verbessert im Nachhinein.
	Außerdem haben viele Lag-Spikes gemeldet, die unterschiedlich stark 
	ausgefallen sind je nach Streckentyp.
	Nach langer Untersuchung konnte die fps-zahl deutlich verbessert werden, 
	jedoch konnte sie nicht maximiert werden. Dies liegt sehr wahrscheinlich 
	an der Zusammenführung der beiden Team-Teile und nach energischer 
	Recherchierung konnten wieder kleiner Verbesserungen stattfinden, jedoch 
	konnte der große Abfall an fps, sobald die race Simulation startet, nicht 
	vollständig behoben werden. Dennoch ist man an einen zufriedenstellenden 
	Punkt gelangt. 
	Weiterhin wurde angemerkt, dass die Energieklau- und Schildfunktion 
	schlecht bis garnicht zu sehen waren. Beide Funktionen wurden 
	besser visualisiert, wobei das Energie-Klauen eine Animation des Fahrers 
	erfordert hätte, welches jedoch aus dem implementierungsplan gestrichen
	werden musste.
	Zusammenfassend fanden alle Propanden das Spiel sehr schön, ereignisreich 
	und erfreulich. Vorallem wurde die Grafik sowie die Steuerung gelobt.
	Weiterhin wurde gesagt, dass man es gut fand, dass Geschwindigkeiten 
	gedrosselt werden konnten durch das Energie klauen und so Chancen Gleichheit
	besteht.
	Es wurde zugesichert, Spaß gehabt zu haben.
	\pagebreak
	\section{Statistiken}

		\subsection{Code-Statistiken}

			\begin{itemize}
				\item Klassen: 134
				\begin{itemize}
					\item Controls: 2
					\item Assets: 16
					\item Console: 1
					\item Profiling: 5
					\item Threading: 2
					\item DummySimulation: 2
					\item Datastructure: 2
					\item Generation: 38
					\item Menu: 24
					\item Rendering: 39
				\end{itemize}
				\item SourceCode: 24.392
				\item SourceCode(von uns erstellt): 23.118
				\item JavaCode: 13.726
				\item JavaCode(von uns erstellt): 13.200
				
			\end{itemize}



		\subsection{Code-Coverage}

		\begin{figure}[htbp]
			\centering
			\includegraphics[width=\linewidth]{./Bilder/coveradeValaris.png}
			\caption{Valaris Code-Coverage}
		\end{figure}


	\pagebreak

	\section{Fazit}

	Abschließend lässt sich sagen, dass das Spiel als ganzes läuft, und beinahe alle Pflichtkriterien erreicht wurden.
	Außerdem lässt sich die Szene durch ein FlyCam betrachten, wobei eine hohe framerate erreicht wird.

	Da uns am Ende die Zeit knapp wurde, konnten wir manche Features nicht implemetieren.
	Zum Beispiel könnte man weitere Assets hinzugefügt werden, um
	den Generierungsalgorithmen noch mehr Raum für Varianz zu geben.
	Sound sowie Splittscrenn fällt auch in diese Kategorie.
	Jedoch haben wir das Meiste erreicht was uns wichtig erschien.

	Im Laufe der Testphase stellte die Intergration uns vor eine große Herausforderung.
	Hier ergaben sich eine Vielzahl an Bugs, manche mehr, manche weniger schlimm.
	Die Integration mit dem anderen Team stellte sich schwieriger heraus als anfangs gedacht.
	Besonders die unterschiedlichen Streckenmodelle und die Umrechnung von einem in das andere führte zu vielen Fehlern und
	war aufwändig zu Debuggen.

	Besonders in den letzten Tagen kamen wir dann aber zu großen Fortschritten
	und konnten die meisten großen Bugs beheben.
	Hinzu konnten wir uns um Feintuning von Parametern
	und das Erstellen und Anpassen von Assets kümmern,
	was dem Spiel Grafisch noch etwas feinschlif gab.
\end{document}